\chapter{Обзор литературы}
\label{ch:chapter1}

Современные рекомендательные системы для маркетплейсов эволюционировали от матричной факторизации к глубоким последовательным моделям, графовым нейросетям и генеративным архитектурам триллионного масштаба. 
Ключевой тренд 2020–2025 годов~--- переход от чисто позиционного моделирования последовательностей к мультимодальному учёту времени, типов действий, поисковых запросов и структуры сессий.
В настоящей главе систематизированы основные направления исследований, формирующие теоретический фундамент для задачи построения рекомендаций с использованием контекстной и сессионной информации пользователя.

\section{Формальная постановка задачи последовательной и сессионной рекомендации}
\label{sec:ch1:problem}

\subsection{Основные определения и нотация}
\label{subsec:ch1:notation}

Пусть $\mathcal{U} = \{u_1, u_2, \ldots, u_N\}$~--- множество пользователей, $\mathcal{I} = \{i_1, i_2, \ldots, i_M\}$~--- множество товаров (каталог маркетплейса). Для каждого пользователя $u \in \mathcal{U}$ определена хронологически упорядоченная последовательность взаимодействий $\mathcal{S}^u = (S^u_1, S^u_2, \ldots, S^u_{|\mathcal{S}^u|})$, где $S^u_t \in \mathcal{I}$~--- товар, с которым пользователь взаимодействовал в момент времени $t$. 
В сессионной постановке анонимная сессия определяется как упорядоченный по времени список $s = [v_{s,1}, v_{s,2}, \ldots, v_{s,n}]$, где $v_{s,i} \in \mathcal{V}$~--- товар, на который пользователь кликнул в рамках сессии $s$ \cite{wu2019srgnn}.

\textbf{Задача предсказания следующего товара} (next-item prediction) формулируется следующим образом: по данной последовательности $(S^u_1, \ldots, S^u_t)$ необходимо определить $S^u_{t+1}$. 
Модель вычисляет релевантности каждого товара $i$ как
\begin{equation}
    r_{i,t} = \mathbf{F}_t^{(b)} \cdot \mathbf{M}_i^\top,
\end{equation}
где $\mathbf{F}_t^{(b)} \in \mathbb{R}^d$~--- выход $b$-го блока модели на позиции $t$, $\mathbf{M}_i \in \mathbb{R}^d$~--- эмбеддинг товара $i$. В сессионной постановке выходной вектор вероятностей определяется как
\begin{equation}
    \hat{\mathbf{y}} = \text{softmax}(\mathbf{s}_h^\top \mathbf{V}),
\end{equation}
где $\mathbf{s}_h$~--- представление сессии, $\mathbf{V}$~--- матрица эмбеддингов всех товаров.

\subsection{Функции потерь}
\label{subsec:ch1:losses}

В литературе используются четыре основных класса функций потерь.

\textbf{Binary Cross-Entropy (BCE)}, применяемая в SASRec \cite{kang2018sasrec}, вычисляется попарно с одним негативным примером на каждую позицию:
\begin{equation}
    \mathcal{L}_{\text{BCE}} = -\sum_{u \in \mathcal{U}} \sum_{t=1}^{|\mathcal{S}^u|} \left[ \log \sigma(r_{o_t, t}) + \log(1 - \sigma(r_{o'_t, t})) \right].
\end{equation}

\textbf{BPR-loss} (Bayesian Personalized Ranking) \cite{rendle2009bpr} оптимизирует попарное ранжирование:
\begin{equation}
    \mathcal{L}_{\text{BPR}} = -\sum_{(u,i,j) \in D_S} \ln \sigma(\hat{x}_{ui} - \hat{x}_{uj}) + \lambda_\Theta \|\Theta\|^2,
\end{equation}
где $D_S = \{(u, i, j) \mid i \in \mathcal{I}^+_u \wedge j \in \mathcal{I} \setminus \mathcal{I}^+_u\}$.

\textbf{Cross-Entropy} по полному каталогу (SR-GNN, BERT4Rec) определяется как $\mathcal{L}_{\text{CE}} = -\log \hat{y}_{v^*}$, а при вычислительной неподъёмности полного softmax применяется \textbf{Sampled Softmax}, ограничивающий знаменатель подмножеством отобранных негативов. 
Отдельно развивается направление \textbf{контрастивных потерь} (InfoNCE), впервые применённых к последовательным рекомендациям в CL4SRec \cite{xie2022cl4srec}.

\subsection{Метрики качества}
\label{subsec:ch1:metrics}

Стандартные метрики оценки включают \textbf{HitRate@K} (доля случаев, когда целевой товар попал в top-$K$), \textbf{MRR@K} (Mean Reciprocal Rank~--- среднее обратного ранга целевого товара), \textbf{NDCG@K} (Normalized Discounted Cumulative Gain) и \textbf{Recall@K}. 
Важно отметить, что Krichene и Rendle \cite{krichene2020sampled} показали, что метрики, вычисленные на подвыборке негативов, могут давать несогласованные рейтинги моделей; в текущей практике предпочтительна оценка по полному каталогу.

\section{Сессионные рекомендательные системы: от RNN к Transformer и далее}
\label{sec:ch1:sequential}

\subsection{Основополагающие архитектуры 2015–2019}
\label{subsec:ch1:foundational}

Первопроходцем в применении глубокого обучения к сессионным рекомендациям стала модель \textbf{GRU4Rec} \cite{hidasi2016gru4rec}. 
Авторы предложили использовать Gated Recurrent Unit для обработки последовательности кликов в анонимной сессии, введя session-parallel mini-batches для эффективного обучения и специализированные ранжирующие функции потерь (BPR и TOP1). 
Формула обновления скрытого состояния: $h_t = \text{GRU}(x_t, h_{t-1})$, где $x_t$~--- эмбеддинг товара на шаге $t$. 
Улучшенная версия с Top-k gains \cite{hidasi2018improved} дала прирост ~35\% относительно оригинала.

\textbf{NARM} \cite{li2017narm} расширила подход RNN гибридным энкодером с механизмом внимания. 
Модель использует два параллельных GRU: глобальный для захвата последовательного поведения и локальный с attention для определения основной цели пользователя в сессии. 
Представление сессии формируется как конкатенация $c_t = [c_g; c_l]$, где $c_g = h_T^g$~--- финальное скрытое состояние глобального энкодера, а $c_l = \sum \alpha_t h_t^l$~--- взвешенная сумма состояний локального.

\textbf{STAMP} \cite{liu2018stamp} отказалась от рекуррентности, заменив её на простой attention-механизм с MLP. 
Модель захватывает общие интересы (средний пулинг эмбеддингов сессии) и текущий интерес (эмбеддинг последнего клика) через трилинейную комбинацию:
\begin{equation}
    \hat{y} = \text{softmax}\!\left(\tanh(W_s m_a) \odot \tanh(W_x x_t) \cdot \mathbf{E}^\top\right).
\end{equation}

\textbf{SASRec} \cite{kang2018sasrec} применила архитектуру Transformer-декодера (однонаправленное каузальное внимание) к последовательной рекомендации, используя механизм внимания на основе масштабированного скалярного произведения с позиционными эмбеддингами. 
Модель достигла на порядок более быстрого обучения по сравнению с CNN/RNN подходами.

\textbf{BERT4Rec} \cite{sun2019bert4rec} адаптировала двунаправленный BERT к задаче через предсказание замаскированного токена, позволив каждому элементу учитывать контекст и слева, и справа. 
Маскирование с вероятностью $\rho \in [0.2, 0.6]$ генерирует больше обучающих сигналов по сравнению с авторегрессионным подходом.

\subsection{Калибровка и эффективные альтернативы 2022–2025}
\label{subsec:ch1:efficient}

Работа \textbf{gSASRec} \cite{petrov2023gsasrec} выявила проблему \textbf{overconfidence} в SASRec при обучении с BCE и негативным сэмплингом: предсказанные вероятности для позитивных товаров систематически завышены, что ухудшает top-$K$ ранжирование. 
Решение~--- обобщённая BCE-потеря (gBCE) с параметром калибровки $\beta$ и увеличенным числом негативов на каждый позитив. 
Архитектура идентична SASRec; меняется только процедура обучения.

\textbf{FMLP-Rec} \cite{zhou2022fmlprec} предложила полностью MLP-архитектуру, заменяющую self-attention обучаемыми фильтрами в частотной области: $\tilde{F}' = \tilde{F} \odot W$ в пространстве FFT, с обратным преобразованием для восстановления. 
Сложность снижается с $O(n^2)$ до $O(n \log n)$ при превосходстве над SASRec и BERT4Rec. 
в работе используется линейный механизм внимания с \textbf{LinRec} \cite{liu2023linrec} используется L2-нормализованное линейное внимание, достигающее $O(N)$ сложности и применяемый как прямая замена стандартного слоя внимания.

Наиболее масштабная работа~--- \textbf{HSTU} \cite{zhai2024hstu} (Meta AI), переформулирующая рекомендации как задачу последовательной трансдукции с использованием Hierarchical Sequential Transduction Units. HSTU заменяет softmax-внимание поэлементной нормализацией (SiLU-гейтинг) и демонстрирует законы масштабирования до \textbf{триллиона параметров} с приростом 12.4\% в онлайн A/B-тестах Meta. 
\textbf{Mamba4Rec} \cite{liu2024mamba4rec} адаптирует архитектуру Mamba (selective state space model) для последовательных рекомендаций с линейной сложностью.

\subsection{Session-aware трансформеры: иерархическое кодирование сессий}
\label{subsec:ch1:session_aware_transformers}

Ключевым направлением развития трансформерных рекомендаторов стало введение явной двухуровневой иерархии сессий, позволяющей одновременно использовать всю историю прошлых сессий пользователя и контекст текущей активной сессии~\cite{quadrana2017hrnn}.

\textbf{HRNN} \cite{quadrana2017hrnn}~--- первая модель, явно разграничивающая внутрисессионную и межсессионную динамику для известного пользователя. 
Архитектура состоит из двух слоёв GRU: сессионного GRU, обрабатывающего объекты внутри сессии, и пользовательского GRU, обновляемого в конце каждой сессии и кодирующего эволюцию предпочтений:
\begin{equation}
    \mathbf{c}_{u,m} = \mathrm{GRU}_\text{usr}\!\left(\mathbf{s}_{u,m},\; \mathbf{c}_{u,m-1}\right),
\end{equation}
где $\mathbf{s}_{u,m}$~--- последнее скрытое состояние сессионного GRU за $m$-ю сессию. 
Эксперименты на данных XING показывают улучшение Recall@5 до $+11\%$ по сравнению с одноуровневыми RNN. Архитектурный шаблон «сессионный энкодер + надсессионный агрегатор» остаётся актуальным и для трансформерных моделей.

\textbf{HCAN} \cite{cui2019hierarchical} переносит шаблон HRNN в трансформерную парадигму через явную двухуровневую архитектуру.
На первом уровне каждая сессия кодируется независимо трансформерным энкодером:
\begin{equation}
    \mathbf{s}_k = \mathrm{Transformer}_\text{item}\!\left(\mathcal{B}_k\right),
\end{equation}
на втором уровне последовательность сессионных представлений $(\mathbf{s}_1, \ldots, \mathbf{s}_{M-1})$ обрабатывается трансформером сессий: $\mathbf{u} = \mathrm{Transformer}_\text{session}(\mathbf{s}_1, \ldots, \mathbf{s}_{M-1})$. 
Оценка следующего объекта формируется через перекрёстное внимание между пользовательским вектором и объектами текущей сессии. 
Чёткое разделение градиентных путей упрощает интерпретацию модели и позволяет обучать уровни независимо.

\textbf{PISA} \cite{tran2024pisa}, разработанный для рекомендования музыки на Deezer, строит представления на уровне сессий, применяя когнитивную модель памяти ACT-R к взвешиванию прошлых сессий. 
Вес $k$-й сессии определяется экспоненциальным затуханием и частотой активации:
\begin{equation}
    w_k \propto \exp(-d \cdot \Delta_k) \cdot f_k^\gamma,
\end{equation}
где $\Delta_k$~--- временной разрыв с момента $k$-й сессии, $f_k$~--- частота объектов, $d$, $\gamma$~--- обучаемые параметры. 
Модель явно учитывает повторные взаимодействия~--- феномен, типичный как для музыки, так и для e-commerce с периодически возвращающимися покупателями.

\textbf{MOJITO} \cite{tran2023mojito} моделирует пользовательские предпочтения как смесь Гауссиан: каждая голова внимания специализируется на одном временно'м режиме (внутрисессионная активность, еженедельный ритм, сезонный тренд), а шлюзовая сеть адаптивно комбинирует их в зависимости от контекста:
\begin{equation}
    \mathrm{MixAttn}(\mathbf{Q},\mathbf{K},\mathbf{V};\mathbf{c}_t)
    = \sum_{k=1}^{K} g_k(\mathbf{c}_t)\;\mathrm{Attn}_k(\mathbf{Q},\mathbf{K},\mathbf{V}),
    \quad \sum_{k} g_k = 1.
\end{equation}
Такая архитектура захватывает одновременно краткосрочную внутрисессионную динамику и долгосрочные сезонные паттерны, что особенно ценно для e-commerce с суточными, еженедельными и сезонными пиками активности.

\textbf{CARCA} \cite{rashed2022carca} вводит перекрёстное внимание между всей историей пользователя и атрибутами объектов-кандидатов:
\begin{equation}
    \hat{r}_{u,i} = \mathrm{CrossAttn}\!\left(\mathbf{H}^{(L)},\; \mathrm{MLP}(\mathbf{a}_i)\right),
\end{equation}
где $\mathbf{H}^{(L)}$~--- вся история пользователя после $L$ трансформерных слоёв, $\mathbf{a}_i$~--- вектор атрибутов кандидата (категория, цена, бренд). 
Это позволяет всем взаимодействиям истории «консультироваться» с атрибутами объекта-кандидата, а не только последнему взаимодействию, что выгодно отличает подход от стандартного скалярного произведения.

\subsection{Самообучение и контрастивные методы для session-aware рекомендаций}
\label{subsec:ch1:ssl_session}

\textbf{S\textsuperscript{3}-Rec} \cite{zhou2020s3rec} предобучает трансформер путём максимизации взаимной информации между четырьмя типами корреляций: атрибут--объект, объект--объект, подпоследовательность--объект, последовательность--объект:
\begin{equation}
    \mathcal{L}_\text{MI}
    = -\mathbb{E}\!\left[
        \log \sigma\!\left(f(e_i, a_i)\right)
        + \log\!\left(1 - \sigma\!\left(f(e_i, a_j)\right)\right)
    \right].
\end{equation}
Предобучение на подпоследовательностях естественно обучает модель различать внутрисессионные и межсессионные паттерны; 
предобучение на атрибутах помогает в сценарии с разреженными историями новых пользователей.

\textbf{DuoRec} \cite{qiu2022duorec} диагностирует коллапс эмбеддингов в трансформерных рекомендаторах и предлагает двойную контрастивную стратегию: метод аугментации данных без учителя (одна история проходит дважды с разными dropout-масками) и метод аугментации данных с учителем (истории, содержащие общий объект, считаются семантически близкими). 
Потеря оптимизируется по NT-Xent:
\begin{equation}
    \mathcal{L}_\text{CL}
    = -\log
    \frac{
      \exp(\mathrm{sim}(\mathbf{z}_i, \mathbf{z}_j^+) / \tau)
    }{
      \sum_{k \neq i} \exp(\mathrm{sim}(\mathbf{z}_i, \mathbf{z}_k) / \tau)
    }.
\end{equation}
Аугментации данных с учителем неявно кодирует общий сессионный контекст: пользователи, взаимодействовавшие с одним объектом в разных сессиях, должны иметь близкие представления.

\textbf{BSARec} \cite{shin2024bsarec} устраняет низкочастотное смещение стандартного самовнимания, добавляя обучаемую высокочастотную коррекцию через преобразование Фурье:
\begin{equation}
    \mathrm{BSA}(\mathbf{H})
    = \mathrm{Attn}(\mathbf{H})
      + \alpha \cdot \mathcal{F}^{-1}\!\left(\mathrm{Rescale}\!\left(\mathcal{F}(\mathbf{H})\right)\right).
\end{equation}
Частотная декомпозиция естественно соответствует двухуровневой природе session-aware задачи: долгосрочные межсессионные предпочтения несут низкочастотный сигнал, краткосрочная внутрисессионная динамика~--- высокочастотный. BSARec без дополнительного контрастивного обучения превосходит DuoRec и FEARec \cite{du2023fearec} на шести стандартных бенчмарках.

\section{Учёт временной информации: от интервалов до ротационных эмбеддингов}
\label{sec:ch1:temporal}

\subsection{Временные интервалы в механизме внимания}
\label{subsec:ch1:time_intervals}

\textbf{TiSASRec} \cite{li2020tisasrec} расширяет SASRec, явно моделируя как абсолютные позиции элементов, так и относительные временные интервалы между взаимодействиями. 
Интервалы между парами элементов $i$ и $j$ обрезаются до максимума $k$: $r_{ij} = \min(|t_i - t_j|, k)$ и отображаются в обучаемые эмбеддинги, модифицирующие attention:
\begin{equation}
    \alpha_{ij} = \frac{q_i \cdot k_j + q_i \cdot r_{ij}^K}{\sqrt{d}}.
\end{equation}
Комбинация абсолютных позиций и относительных временных интервалов превосходит использование каждого компонента по отдельности на разреженных и плотных датасетах.

\subsection{Периодические кодирования времени}
\label{subsec:ch1:time_encoding}

\textbf{Time2Vec} \cite{kazemi2020time2vec} предложила универсальное модельно-агностичное векторное представление времени, объединяющее линейный тренд и $k$ периодических синусоидальных компонент:
\begin{equation}
    \text{t2v}(\tau)[i] = \begin{cases}
        \omega_i \tau + \varphi_i, & i = 0, \\
        \sin(\omega_i \tau + \varphi_i), & 1 \leq i \leq k,
    \end{cases}
\end{equation}
где $\omega_i, \varphi_i$~--- обучаемые частоты и фазовые сдвиги.

\subsection{Ротационные позиционные эмбеддинги в рекомендациях}
\label{subsec:ch1:rope}

Оригинальный \textbf{RoFormer} \cite{su2021roformer} кодирует позицию $m$ через матрицу вращения $R_m = \text{diag}(R_1(m\theta_1), \ldots)$, так что скалярное произведение $\langle R_m q, R_n k \rangle$ зависит только от относительной позиции $m - n$.

\textbf{TO-RoPE} \cite{wei2025torope} адаптирует RoPE для рекомендаций, совместно кодируя дискретный индекс и непрерывное время через три стратегии слияния: early fusion (совмещение углов в одной плоскости), split-by-dim (отдельные плоскости для индекса и времени) и split-by-head (выделение целых голов под каждый тип). 
\textbf{RecGRELA} \cite{hu2025recgrela} комбинирует RoPE с линейным attention и SiLU-гейтом для $O(N)$ сложности, устраняя необходимость в обучаемых таблицах позиционных эмбеддингов.

\subsection{Устаревание эмбеддингов и временной спад}
\label{subsec:ch1:staleness}

Проблема staleness (устаревания) эмбеддингов особенно актуальна для маркетплейсов с быстрым оборотом ассортимента. 
HSTU \cite{zhai2024hstu} использует логарифмическую бакетизацию разностей временных меток:
\begin{equation}
    \text{bucket}(a_{ij}) = \left\lfloor \frac{\log(\max(1, |a_{ij}|))}{0.301} \right\rfloor,
\end{equation}
отображая каждый бакет в обучаемый скалярный bias в матрице внимания. 
Работа \cite{meta2025semantic} показала, что при случайном хешировании \textbf{половина товаров покидает каталог в течение 6 дней}, делая эмбеддинги нестабильными; 
семантические идентификаторы на основе контентных признаков обеспечивают значительно более стабильные представления. 
Классический фундамент заложен Koren \cite{koren2009timeSVD} в работе TimeSVD++, где латентные факторы пользователей и товаров эволюционируют как функции времени.

\section{Гибридные эмбеддинги и проблема dimensional collapse}
\label{sec:ch1:embeddings}

\subsection{Deep \& Cross Network для формирования представлений}
\label{subsec:ch1:dcn}

\textbf{DCN} \cite{wang2017dcn} заменил ручное конструирование перекрёстных признаков алгоритмическим cross-network, где каждый слой вычисляет:
\begin{equation}
    \mathbf{x}_{l+1} = \mathbf{x}_0 \cdot \mathbf{x}_l^\top \cdot \mathbf{w}_l + \mathbf{b}_l + \mathbf{x}_l.
\end{equation}

\textbf{DCN V2} \cite{wang2021dcnv2} заменил весовой вектор полной матрицей $W_l \in \mathbb{R}^{d \times d}$:
\begin{equation}
    \mathbf{x}_{l+1} = \mathbf{x}_0 \odot (W_l \cdot \mathbf{x}_l + \mathbf{b}_l) + \mathbf{x}_l.
\end{equation}
Эмпирически обнаруженная \textbf{низкоранговая структура} обученных матриц мотивировала вариант DCN-Mix с mixture-of-low-rank-experts, достигающий сопоставимой точности при \textasciitilde30\% снижении параметров/FLOPs. \textbf{GDCN} \cite{wang2023gdcn} добавил информационные gates поверх каждого cross-слоя для селективной фильтрации шумных высокопорядковых взаимодействий. Другие ключевые гибридные архитектуры включают \textbf{Wide \& Deep} \cite{cheng2016widedeep} (Google Play), \textbf{NCF} \cite{he2017ncf} и \textbf{DLRM} \cite{naumov2019dlrm} (Meta).

\subsection{Dimensional collapse в рекомендательных моделях}
\label{subsec:ch1:collapse}

Работа \cite{guo2024collapse} идентифицировала \textbf{embedding collapse}~--- тенденцию эмбеддинговых матриц занимать низкоразмерное подпространство даже при увеличении размерности. 
Диагностика проводится через SVD-разложение $E = U\Sigma V^\top$; метрика Information Abundance $= \sum \sigma_i / \sigma_{\max}$ количественно оценивает степень коллапса. 
Авторы предложили архитектуру с несколькими эмбеддингами и специализированными модулями для различных типов взаимодействий.Chen et al. \cite{chen2024noncl} показали, что контрастивные GNN-модели также подвержены коллапсу размерности, и предложили подход неконтрастивного обучения (nCL) в гиперсферическом пространстве.

\subsection{Hash-based эмбеддинги для больших каталогов}
\label{subsec:ch1:hash}

Для каталогов с миллионами товаров полная эмбеддинговая таблица $O(M \times d)$ становится узким местом по памяти. \textbf{Compositional Embeddings} \cite{shi2020compositional} используют $k$ малых таблиц с комплементарными разбиениями:
\begin{equation}
    e(c) = \text{OP}(T_1[h_1(c)], \ldots, T_k[h_k(c)]),
\end{equation}
сжимая память до $O(k \cdot |S|^{1/k} \times D)$~--- до \textbf{15$\times$ сжатие} при незначительной потере точности на Criteo/DLRM. 
\textbf{Deep Hash Embedding} \cite{kang2020deepHash} полностью заменяет таблицы нейросетью: $e(x) = \text{DNN}(\text{transform}(h_1(x), \ldots, h_k(x)))$, гарантируя нулевые коллизии хешей при 75\% сжатии параметров. 
\textbf{Frequency-Based Double Hashing} \cite{zhang2020doublehash} использует двойное хеширование с частотно-зависимым распределением ёмкости.

\section{Контекстуальная информация: запросы, типы действий, сессионные токены}
\label{sec:ch1:context}

\subsection{Учёт поисковых запросов в последовательных рекомендациях}
\label{subsec:ch1:queries}

\textbf{Query-Aware Sequential Recommendation} \cite{he2022queryaware}~--- первая работа, формализующая совместное моделирование поисковых запросов и последовательностей товаров. Модель организует запросы и товары как гетерогенные последовательности и применяет self-attention по объединённому представлению; разные запросы над одной последовательностью дают разные рекомендации.

\textbf{JSR} \cite{zamani2020jsr} совместно обучает модели поиска и рекомендаций через общие матрицы эмбеддингов $W$ (пользователи) и $E$ (товары), оптимизируя объединённую функцию потерь $L = L_{IR} + L_{RS}$.

\textbf{UDITSR} \cite{uditsr2024} моделирует двойные интенты (явные из поиска, неявные из рекомендаций) через граф двух сценариев с контрастивной потерей.

\textbf{SynerGen} \cite{synergen2025}~--- первая decoder-only генеративная модель, поддерживающая и query-aware поиск, и query-free рекомендации в едином backbone с time-aware rotary positional embedding.

\subsection{Разделение типов действий (multi-behavior recommendation)}
\label{subsec:ch1:multibehavior}

В e-commerce пользователь совершает различные действия: клик, добавление в корзину, покупка, просмотр. \textbf{MB-STR} \cite{yuan2022mbstr} использует Transformer с behavior-type-specific routing через различные feed-forward слои и sparse Mixture-of-Experts для кодирования гетерогенных последовательных паттернов. Обучение ведётся через behavior-aware masked item prediction: товары маскируются, а типы действий остаются видимыми.

\textbf{NMTR} \cite{gao2019nmtr} формализует каскадные отношения между поведениями ($\text{view} \to \text{cart} \to \text{purchase}$) как связанные задачи multi-task обучения. \textbf{MBSRec} \cite{mbsrec2024} захватывает зависимости между поведениями через multi-head self-attention с взвешенной BCE-потерей. Обзор Kim et al. \cite{kim2026survey} систематизирует всю область в трёхэтапный фреймворк: Data Modeling, Encoding, Training.

\subsection{Негативная выборка из просмотренных товаров}
\label{subsec:ch1:negative}

Ding et al. \cite{ding2018view} в работе «Improving Implicit Recommender Systems with View Data» установили \textbf{трёхуровневую иерархию предпочтений}: $\text{purchased} > \text{viewed} > \text{non-viewed}$. Просмотренные, но некупленные товары служат информативными hard negatives~--- сложнее случайных, но мягче истинных негативов.

Shi et al. \cite{shi2023hardneg} впервые дали теоретическое обоснование hard negative sampling, доказав, что BPR с dynamic negative sampling является точным оценщиком OPAUC. Prakash et al. \cite{prakash2024neg} оценили стратегии негативного сэмплирования специфично для SASRec-архитектуры, выявив, что адаптивная смешанная стратегия достигает наилучших результатов.

\subsection{Сессионные разделительные токены}
\label{subsec:ch1:session_tokens}

Seol, Ko, Lee \cite{seol2022session} в работе «Exploiting Session Information in BERT-based Session-aware Sequential Recommendation» предложили три механизма включения сессионной информации в BERT4Rec:
\begin{enumerate}
    \item \textbf{Session Token (ST)}~--- обучаемый специальный токен, вставляемый между сессиями (аналог [SEP] в NLP);
    \item \textbf{Session Segment Embedding (SSE)}~--- обучаемый эмбеддинг уровня сессии, добавляемый к каждому элементу: $\mathbf{x} = \mathbf{IE}_i + \mathbf{PE}_p + \mathbf{SSE}_j$;
    \item \textbf{Time-Aware Self-Attention (TAS)}~--- временное кодирование промежутков между взаимодействиями.
\end{enumerate}
Все три подхода улучшают vanilla BERT4Rec/SASRec без сложных иерархических архитектур. \textbf{Transformers4Rec} \cite{moreira2021transformers4rec} (NVIDIA) предоставляет инфраструктуру для применения NLP-концепций (включая специальные токены) к рекомендациям, выиграв два соревнования по session-based рекомендациям.

\subsection{Моделирование эволюции интересов пользователя}
\label{subsec:ch1:interest_evolution}

Интент пользователя не является статичным: он претерпевает непрерывную эволюцию как в пределах одной сессии~--- от этапа формирования потребности к целенаправленному поиску~--- так и на межсессионном горизонте, под влиянием сезонных факторов и изменений жизненного контекста. Данное явление получило наименование \textit{interest drift} и составляет самостоятельное направление исследований.

\textbf{DIN} (Deep Interest Network)~\cite{zhou2018din} ввёл понятие \textit{адаптивного активационного механизма}: вместо агрегации всей пользовательской истории в единый вектор модель вычисляет вес релевантности каждого исторического взаимодействия относительно конкретного товара-кандидата $\mathbf{v}_a$:
\begin{equation}
    \mathbf{e}_u = \sum_{i \in \mathcal{H}^u} a(\mathbf{v}_i,\, \mathbf{v}_a) \cdot \mathbf{v}_i,
    \qquad
    a(\mathbf{v}_i, \mathbf{v}_a) = \sigma\!\left(\mathrm{MLP}([\mathbf{v}_i;\, \mathbf{v}_a;\, \mathbf{v}_i \odot \mathbf{v}_a])\right).
\end{equation}
Тем самым формируется \textit{инстанс-зависимое} представление пользователя: его эффективный профиль адаптируется к каждому кандидату независимо, что принципиально отличается от инстанс-независимых методов усредняющего пулинга.

\textbf{DIEN} (Deep Interest Evolution Network)~\cite{zhou2019dien} расширяет DIN введением явного механизма эволюции интересов. Двухуровневая архитектура объединяет \textit{слой извлечения интересов} (GRU) и \textit{слой эволюции интересов} (AUGRU~--- GRU с внимания-взвешенным update gate), позволяя отслеживать только ту часть траектории интересов, которая релевантна текущему кандидату:
\begin{equation}
    \mathbf{h}'_t = (1 - \mathbf{u}'_t) \odot \mathbf{h}'_{t-1} + \mathbf{u}'_t \odot \tilde{\mathbf{h}}'_t,
    \qquad
    \mathbf{u}'_t = a(\mathbf{h}_t,\, \mathbf{v}_a) \cdot \mathbf{u}_t,
\end{equation}
где $\mathbf{u}_t$~--- стандартный update gate GRU, масштабируемый скором внимания к текущему кандидату. В контексте session-aware задачи DIEN формализует ключевую гипотезу: не все прошлые сессии одинаково релевантны текущему намерению пользователя.

\textbf{SIM} (Search-based Interest Model)~\cite{pi2020sim} решает задачу масштабируемого моделирования долгосрочной истории для пользователей с $10^4$ и более взаимодействиями. На первом этапе из всей истории жёстким или мягким поиском (GSU/ESU) извлекаются топ-$K$ наиболее схожих с кандидатом взаимодействий; на втором к ним применяется механизм внимания с временны'м кодированием. Промышленное развёртывание в системе Alibaba подтвердило значимый прирост качества при масштабировании истории без увеличения вычислительной сложности инференса.

\subsection{Диsentanglement интентов: разделение устойчивых и волатильных компонент}
\label{subsec:ch1:disentangle}

Фундаментальная проблема моделирования интента~--- декомпозиция пользовательских предпочтений на \textit{инвариантную} составляющую (устойчивые долгосрочные интересы, не зависящие от текущего контекста) и \textit{контекстно-зависимую} составляющую (волатильные краткосрочные намерения, специфичные для текущей сессии). Методы без явного разделения систематически смешивают эти сигналы, что приводит к ошибкам интерпретации.

\textbf{CauseRec}~\cite{zhang2021causerec} применяет инструментарий каузального вывода для диentanglement представления пользователя на \textit{стабильный интерес} $\mathbf{z}_s$ и \textit{нестабильный интерес} $\mathbf{z}_u$. Контрфактуальное обучение с интервенциями на историю взаимодействий обеспечивает статистическую независимость компонент $\mathbf{z}_s \perp\!\!\!\perp \mathbf{z}_u$ и устойчивость стабильного интереса к случайным возмущениям сессионного контекста.

\textbf{ICLRec} (Intent Contrastive Learning for Sequential Recommendation)~\cite{chen2022iclrec} моделирует множество латентных интентов через прототипическое обучение: $K$ кластерных центроидов $\{\mathbf{c}_k\}$ в пространстве пользовательских эмбеддингов представляют дискретные прототипы намерений. Контрастивная цель максимизирует согласованность между представлением пользователя и его ближайшим прототипом:
\begin{equation}
    \mathcal{L}_\text{intent}
    = -\log \frac{
        \exp(\mathbf{z}_u \cdot \mathbf{c}_{k^*} / \tau)
    }{
        \sum_{j=1}^{K} \exp(\mathbf{z}_u \cdot \mathbf{c}_j / \tau)
    }, \qquad k^* = \operatorname*{argmax}_k\, \mathbf{z}_u \cdot \mathbf{c}_k.
\end{equation}
Совместное обучение с основной задачей предсказания следующего товара обеспечивает семантическую интерпретируемость кластеров и прирост NDCG@10 на 2.5--5\% на бенчмарках Amazon и Yelp.

\textbf{CORE} (Consistent Representation Learning for Session-based Recommendation)~\cite{hou2022core} постулирует требование \textit{семантической согласованности}: представление сессии и представления составляющих её товаров должны принадлежать единому семантическому пространству. Это ограничение гарантирует, что сессионный вектор остаётся интерпретируемым в терминах товарных эмбеддингов, а не является произвольной скрытой переменной другого пространства.

\subsection{Обогащение контекста через внешние структурированные знания}
\label{subsec:ch1:kgraph}

Идентификаторные модели (ID-based) ограничены в способности к семантическому обобщению: близость двух товаров определяется исключительно статистикой совместных появлений в историях пользователей, тогда как содержательная семантическая близость~--- принадлежность к одной категории, схожие технические характеристики, единый производитель~--- остаётся скрытой. Для преодоления этого ограничения применяется обогащение признаковых пространств через графы знаний, категориальные иерархии и атрибутивные схемы.

\textbf{KGAT} (Knowledge Graph Attention Network)~\cite{wang2019kgat} рассматривает взаимодействия пользователей и товаров как специальный тип отношений в гетерогенном графе знаний $\mathcal{G} = \{(h, r, t)\}$. Рёберно-специфичный механизм внимания взвешивает соседей по типу отношения:
\begin{equation}
    \pi(h, r, t) = (\mathbf{W}_r \mathbf{e}_t)^\top \tanh(\mathbf{W}_r \mathbf{e}_h + \mathbf{e}_r),
\end{equation}
где $\mathbf{W}_r$~--- матрица проекции, специфичная для типа отношения $r$. Послойное распространение по графу формирует обогащённые эмбеддинги, отражающие семантику категориальных связей, атрибутов и производителей. В session-aware контексте это позволяет устанавливать семантическое родство между товарами, посещёнными в разных сессиях.

\textbf{KGIN} (Knowledge Graph-based Intent Network)~\cite{wang2021kgin} вводит понятие \textit{интент-пути}: цепочка отношений $(r_1, r_2, \ldots, r_L)$ от пользователя до товара интерпретируется как атомарное намерение. Агрегация по всем таким путям порождает интент-специфичное представление пользователя, обеспечивающее интерпретируемость рекомендаций через структуру намерений:
\begin{equation}
    \mathbf{p}_{(r_1,\ldots,r_L)} = \mathbf{e}_{r_1} \odot \mathbf{e}_{r_2} \odot \cdots \odot \mathbf{e}_{r_L},
    \qquad
    \mathbf{u}_\text{intent} = \sum_{\text{path}} \beta_{\text{path}} \cdot \mathbf{p}_{\text{path}}.
\end{equation}
Подход открывает перспективу использования онтологий товарных каталогов (иерархии категорий, атрибутивные схемы маркетплейса) как источника структурированного контекста для уточнения пользовательского интента на уровне сессии.

\section{Графовые подходы к сессионным рекомендациям}
\label{sec:ch1:graph}

\subsection{От SR-GNN к глобальному контексту}
\label{subsec:ch1:srgnn}

\textbf{SR-GNN} \cite{wu2019srgnn}~--- семинальная работа, впервые моделирующая сессионные последовательности как ориентированные графы. Из сессии $s = [v_{s,1}, \ldots, v_{s,n}]$ строится граф $\mathcal{G}_s = (\mathcal{V}_s, \mathcal{E}_s)$, где рёбра $(v_{s,i-1}, v_{s,i})$ нормализуются по исходящей степени. Gated Graph Neural Network обновляет эмбеддинги через механизм гейтов:
\begin{equation}
    \mathbf{v}_i^l = (1 - \mathbf{z}_t) \odot \mathbf{v}_i^{l-1} + \mathbf{z}_t \odot \tilde{\mathbf{v}}_i^l.
\end{equation}
Представление сессии $\mathbf{s}_h = W_3[\mathbf{s}_l; \mathbf{s}_g]$ объединяет локальный (последний клик) и глобальный (attention-взвешенная сумма) компоненты.

\textbf{GCE-GNN} \cite{wang2020gcegnn} устранила ограничение моделирования только внутрисессионных переходов, введя \textbf{два уровня эмбеддингов}: session-level (через GGNN по графу сессии) и global-level (через session-aware attention по глобальному графу всех сессий). Агрегация двух уровней выполняется через soft attention.

\textbf{TAGNN} \cite{yu2020tagnn} добавила target-aware attention, где представление сессии адаптивно меняется в зависимости от целевого товара-кандидата: $\mathbf{s}_t = \sum \alpha_i \cdot \mathbf{v}_i$, где $\alpha_i = \text{softmax}(f(\mathbf{v}_i, \mathbf{v}_t))$.

\textbf{GC-SAN} \cite{xu2019gcsan} комбинирует GNN (для локальных зависимостей) с multi-layer self-attention (для дальних зависимостей), финальное представление~--- взвешенная комбинация $\mathbf{s}_h = \omega \cdot \mathbf{s}_a + (1 - \omega) \cdot \mathbf{v}_n$.

\subsection{Графовая коллаборативная фильтрация}
\label{subsec:ch1:gcf}

\textbf{NGCF} \cite{wang2019ngcf} предложил embedding propagation на двудольном графе «пользователь–товар» для кодирования коллаборативного сигнала высокого порядка:
\begin{equation}
    m_{u \leftarrow i} = \frac{1}{\sqrt{|N_u|}\sqrt{|N_i|}}\left(W_1 \mathbf{e}_i + W_2(\mathbf{e}_i \odot \mathbf{e}_u)\right).
\end{equation}

\textbf{LightGCN} \cite{he2020lightgcn} через ablation-анализ NGCF обнаружил, что feature transformation и нелинейная активация не только не помогают, но вредят в задаче CF. LightGCN сохраняет только \textbf{линейное распространение с агрегацией соседей} и взвешенную сумму по слоям:
\begin{equation}
    \mathbf{e}_u = \sum_{k=0}^{K} \alpha_k \cdot \mathbf{e}_u^{(k)},\quad \alpha_k = \frac{1}{K+1},
\end{equation}
давая \textasciitilde16\% прирост над NGCF.

\subsection{Гетерогенные графы}
\label{subsec:ch1:heterogeneous}

\textbf{HG-GNN} \cite{pang2022hggnn} строит \textbf{гетерогенный глобальный граф} с двумя типами узлов (пользователи, товары) и тремя типами рёбер (исторические взаимодействия, внутрисессионные переходы, глобальные ко-вхождения). Meta-path-based агрегация обучает Current Preference Encoder и Historical Preference Encoder, объединяемые в Personalized Session Encoder. Фундамент гетерогенных графовых сетей заложен в \textbf{HAN} \cite{wang2019han} с иерархическим attention: node-level (важность соседей по мета-пути) и semantic-level (важность различных мета-путей).

\section{Архитектурные улучшения для промышленных систем}
\label{sec:ch1:industrial}

\subsection{Mixture of Experts в рекомендациях}
\label{subsec:ch1:moe}

\textbf{MMoE} \cite{ma2018mmoe} заменил shared-bottom архитектуру multi-task обучения общими экспертами с task-specific gating:
\begin{equation}
    y_k = h_k\!\left(\sum_i g_{k,i}(x) \cdot f_i(x)\right),\quad g_k(x) = \text{softmax}(W_{g_k} \cdot x).
\end{equation}

\textbf{PLE} \cite{tang2020ple} расширил MMoE совместными и task-specific экспертами с прогрессивной послойной экстракцией. Прорыв в балансировке нагрузки экспертов достигнут в работе DeepSeek \cite{wang2024deepseek}: к скорам маршрутизации добавляется обучаемый bias $b_i$ для каждого эксперта, корректируемый по загрузке (если эксперт перегружен~--- $b_i$ уменьшается на $\gamma$, иначе увеличивается). Критически, bias влияет только на маршрутизацию, не на gate-значения, обеспечивая нулевые интерференционные градиенты. Метод развёрнут в \textbf{DeepSeek-V3} (671B параметров, 37B активных).

\subsection{TorchRec и jagged tensors}
\label{subsec:ch1:torchrec}

\textbf{TorchRec} \cite{ivchenko2022torchrec} (Meta)~--- PyTorch-библиотека с примитивами разреженности и параллелизма для рекомендаций промышленного масштаба. Ключевая структура данных~--- \textbf{KeyedJaggedTensor (KJT)}: эффективное хранение категориальных признаков переменной длины через непрерывный массив значений с массивом смещений, исключающее потери на padding. Поддерживает гибридный data+model параллелизм с table-wise, row-wise, column-wise шардированием эмбеддинговых таблиц. Используется для обучения моделей \textbf{более 3 триллионов параметров} в Meta. \textbf{Jagged Flash Attention} обеспечивает 9$\times$ ускорения и 22$\times$ сжатия памяти по сравнению с dense attention.

\subsection{Knowledge distillation и LoRA}
\label{subsec:ch1:distillation}

Knowledge distillation в рекомендациях включает Ranking Distillation \cite{tang2018ranking}, DE-RRD \cite{kang2020derr}, Topology Distillation \cite{kang2021topology} и Privileged Features Distillation \cite{xu2020privileged}~--- передачу знаний из признаков, доступных при обучении, но недоступных при сервинге.

\textbf{LoRA} \cite{hu2022lora} вводит низкоранговую адаптацию $W = W_0 + BA$, где $B \in \mathbb{R}^{d \times r}$, $A \in \mathbb{R}^{r \times k}$, $r \ll \min(d, k)$, сокращая обучаемые параметры в 10000$\times$ для GPT-3. Принцип низкоранговости перекликается с DCN V2 (низкоранговые cross-слои), compositional embeddings (факторизация таблиц) и проблемой dimensional collapse. Применение LoRA к рекомендациям активно развивается в контексте федеративных рекомендательных систем \cite{nguyen2024fedlora} и fine-tuning LLM-based рекомендателей.

\section{Наборы данных для оценки}
\label{sec:ch1:datasets}

Для экспериментальной оценки предлагаемого метода выбраны четыре публично доступных набора данных, охватывающих различные сценарии e-commerce рекомендования: стандартные бенчмарки последовательной рекомендации, мультиповеденческие данные и датасеты с поисковыми запросами.

\subsection{Amazon Reviews}
\label{subsec:ch1:amazon}

Семейство наборов данных \textbf{Amazon Reviews} \cite{mcauley2015amazon} содержит взаимодействия пользователей с товарами нескольких категорий интернет-магазина Amazon. Данные включают оценки, текстовые отзывы, временны'е метки и богатые метаданные товаров (категория, цена, бренд, текстовое описание). Для экспериментов используются две категории с контрастными характеристиками разреженности:

\begin{itemize}
    \item \textbf{Amazon-Beauty}~--- разреженная категория ($\sim$22k пользователей, $\sim$12k товаров), типична для cold-start сценариев;
    \item \textbf{Amazon-Sports}~--- умеренная плотность ($\sim$36k пользователей, $\sim$18k товаров), стандартный бенчмарк для сравнения с SASRec, BERT4Rec, DuoRec, BSARec.
\end{itemize}

Сессии конструируются из хронологических последовательностей взаимодействий путём разбивки по временному порогу (разрыв $>30$ минут~--- новая сессия). Для оценки применяется протокол leave-one-out: последнее взаимодействие каждого пользователя отводится под тест, предпоследнее~--- под валидацию.

\subsection{MovieLens-1M}
\label{subsec:ch1:movielens}

\textbf{MovieLens-1M} \cite{harper2015movielens}~--- классический бенчмарк рекомендательных систем, содержащий около 1 миллиона оценок 6 040 пользователей по 3 706 фильмам с временны'ми метками. Несмотря на специфику предметной области (кино), датасет широко используется для воспроизводимого сравнения с иерархическими моделями: HRNN \cite{quadrana2017hrnn} и Session-aware BERT4Rec \cite{seol2022session} первично оценивались на этом наборе. Плотность данных ($\sim$4.5\% заполненности матрицы взаимодействий) позволяет проверить поведение модели при достаточном количестве обучающих сигналов. Сессии разбиваются аналогично Amazon по временному порогу.

\subsection{Tmall (IJCAI-15)}
\label{subsec:ch1:tmall}

\textbf{Tmall}~--- набор данных крупного китайского маркетплейса Alibaba Tmall, предоставленный для конкурса IJCAI-15. Содержит журналы поведения авторизованных пользователей с четырьмя типами действий: клик, добавление в избранное, добавление в корзину, покупка. Наличие полного спектра поведенческих типов делает датасет незаменимым для оценки multi-behavior компонента. Ключевые характеристики: $\sim$884k пользователей, $\sim$1.1M товаров, $\sim$54M событий. Временны'е метки позволяют восстановить структуру сессий; известные идентификаторы пользователей обеспечивают полную применимость session-aware постановки.

\subsection{Coveo SIGIR eCommerce Dataset}
\label{subsec:ch1:coveo}

\textbf{Coveo SIGIR eCommerce Dataset} \cite{tagliabue2021coveo}~--- единственный из рассматриваемых наборов, содержащий как поисковые запросы, так и поведенческие события. Датасет опубликован в рамках задачи SIGIR 2021 eCom Data Challenge и включает реальные данные e-commerce платформы: поисковые запросы пользователей, просмотры страниц товаров, добавления в корзину и покупки в рамках явно размеченных сессий.

Основные характеристики:
\begin{itemize}
    \item Явные границы сессий (session ID присвоены в исходных данных);
    \item Поисковые запросы на естественном языке, ассоциированные с событиями;
    \item Несколько типов событий (pageview, detail, add, purchase);
    \item Метаданные товаров: категория, цена, описание.
\end{itemize}

Этот датасет является основным для оценки query-aware компонента предлагаемой архитектуры, поскольку позволяет верифицировать, насколько явное включение поисковых запросов улучшает качество рекомендаций по сравнению с моделями, опирающимися только на поведенческую историю.

\subsection{Сравнительная характеристика датасетов}

\begin{table}[h]
\centering
\caption{Сравнение наборов данных для экспериментальной оценки.}
\label{tab:ch1:datasets}
\small
\begin{tabularx}{\textwidth}{@{} l r r r c c c @{}}
\toprule
\textbf{Датасет} & \textbf{Пользователи} & \textbf{Товары} & \textbf{События} &
\textbf{Сессии} & \textbf{Multi-behavior} & \textbf{Запросы} \\
\midrule
Amazon-Beauty  & 22\,363  & 12\,101  & 198\,502  & конструируются & -- & -- \\
Amazon-Sports  & 35\,598  & 18\,357  & 296\,337  & конструируются & -- & -- \\
MovieLens-1M   & 6\,040   & 3\,706   & 1\,000\,209 & конструируются & -- & -- \\
Tmall          & 884\,531 & 1\,090\,390 & 54\,925\,331 & конструируются & \checkmark & -- \\
Coveo SIGIR    & $\sim$1M & $\sim$1.6M & $\sim$50M & явные & \checkmark & \checkmark \\
\bottomrule
\end{tabularx}
\end{table}

Для всех датасетов применяется единый протокол предобработки: удаление пользователей с менее чем 5 взаимодействиями (стандартный порог 5-core), нормализация временны'х меток, разбивка 80/10/10 для трейна, валидации и теста по времени.

\section{Метрики оценки качества рекомендаций}
\label{sec:ch1:metrics_eval}

Для количественной оценки качества рекомендательных моделей в задаче предсказания следующего товара используются метрики ранжирования, вычисляемые по полному каталогу товаров~\cite{krichene2020sampled}. Оценка по подвыборке негативов намеренно исключается, поскольку приводит к несогласованным рейтингам моделей~\cite{krichene2020sampled}.

\subsection{Основные метрики ранжирования}
\label{subsec:ch1:ranking_metrics}

\textbf{Hit Rate@K (HR@K)} измеряет долю тестовых запросов, для которых целевой товар попал в первые $K$ позиций рекомендованного списка:
\begin{equation}
    \mathrm{HR@K} = \frac{1}{|\mathcal{U}|} \sum_{u \in \mathcal{U}} \mathbf{1}\!\left[\mathrm{rank}(i^*_u) \leq K\right],
\end{equation}
где $i^*_u$~--- целевой товар пользователя $u$, $\mathrm{rank}(i^*_u)$~--- его позиция в ранжированном списке. При единственном целевом товаре (leave-one-out) HR@K совпадает с Recall@K.

\textbf{Normalized Discounted Cumulative Gain@K (NDCG@K)} учитывает позицию целевого товара в списке через логарифмическое дисконтирование:
\begin{equation}
    \mathrm{NDCG@K} = \frac{1}{|\mathcal{U}|} \sum_{u \in \mathcal{U}} \frac{\mathbf{1}\!\left[\mathrm{rank}(i^*_u) \leq K\right]}{\log_2\!\left(\mathrm{rank}(i^*_u) + 1\right)}.
\end{equation}
Рекомендация на первой позиции даёт максимальный вклад $1/\log_2(2) = 1$; рекомендация на позиции $K$~--- минимальный в пределах top-$K$.

\textbf{Mean Reciprocal Rank (MRR)} усредняет обратные ранги целевых товаров по всем пользователям:
\begin{equation}
    \mathrm{MRR} = \frac{1}{|\mathcal{U}|} \sum_{u \in \mathcal{U}} \frac{1}{\mathrm{rank}(i^*_u)}.
\end{equation}
MRR особенно чувствителен к точным попаданиям на первых позициях и дополняет NDCG при анализе поведения модели в верхней части рейтинга.

В экспериментах используются значения $K \in \{10, 20\}$ как наиболее распространённые в литературе по session-aware рекомендованию~\cite{kang2018sasrec, sun2019bert4rec, qiu2022duorec, shin2024bsarec}.

\subsection{Дополнительные метрики для session-aware анализа}
\label{subsec:ch1:extended_metrics}

Помимо стандартных метрик ранжирования, для углублённого анализа предлагается вычислять следующие показатели.

\textbf{Метрики в разрезе длины истории.} Пользователи разбиваются на группы по числу прошлых сессий: короткая история ($\leq 3$ сессий), средняя ($4$--$10$) и длинная ($> 10$). Раздельный расчёт HR@10 и NDCG@10 по группам позволяет оценить способность модели работать как в cold-start сценарии, так и при богатой долгосрочной истории.

\textbf{Метрики по типу поведения (для Tmall).} На датасете Tmall целевым событием является покупка; метрики вычисляются отдельно по событиям типа «покупка», что отражает бизнес-релевантную задачу предсказания конверсии.

\textbf{Метрики для query-aware сценария (для Coveo).} На датасете Coveo дополнительно вычисляется HR@10 в подмножестве сессий, начинающихся с поискового запроса, против сессий без запроса. Это позволяет количественно оценить вклад query-aware компонента в качество рекомендаций.

\subsection{Протокол оценки}
\label{subsec:ch1:eval_protocol}

Во всех экспериментах применяется единый протокол:

\begin{enumerate}
    \item \textbf{Разбивка по времени}: обучающая, валидационная и тестовая выборки формируются хронологически в соотношении 80/10/10, исключая утечку данных из будущего.
    \item \textbf{Leave-last-out}: последнее взаимодействие каждого пользователя отводится под тест, предпоследнее~--- под валидацию.
    \item \textbf{Полный каталог}: метрики вычисляются по всем товарам каталога без отбора негативов.
    \item \textbf{5-core фильтрация}: удаляются пользователи и товары с менее чем 5 взаимодействиями.
    \item \textbf{Повторяемость}: для каждого эксперимента фиксируются случайные зёрна (random seed); результаты усредняются по 3 запускам.
\end{enumerate}

\section{Выводы по главе}
\label{sec:ch1:summary}

Проведённый анализ литературы выявляет несколько ключевых наблюдений. Во-первых, ни одна существующая работа не объединяет все четыре типа контекстных сигналов~--- \textbf{поисковые запросы, типы действий, сессионные границы и временную информацию}~--- в единой архитектуре. Каждый из этих сигналов исследован отдельно \cite{he2022queryaware, yuan2022mbstr, seol2022session, li2020tisasrec}, но их синергия остаётся неисследованной.

Во-вторых, применение RoPE к рекомендациям находится на ранней стадии \cite{wei2025torope, hu2025recgrela}, и его комбинация с сессионной структурой и multi-behavior моделированием представляет новую исследовательскую возможность.

В-третьих, проблема dimensional collapse при масштабировании моделей \cite{guo2024collapse} и embedding staleness при быстром обороте каталога \cite{meta2025semantic} требуют архитектурных решений, специфичных для e-commerce с его динамичным ассортиментом.

Наконец, приёмы из LLM-мира~--- \textbf{MoE с auxiliary-loss-free балансировкой} \cite{wang2024deepseek}, \textbf{LoRA для инкрементального обучения} \cite{hu2022lora}, \textbf{flash attention для jagged данных} \cite{ivchenko2022torchrec}~--- создают инфраструктуру для моделей следующего поколения, способных обрабатывать длинные мультимодальные сессии в реальном времени.

Таким образом, данная диссертация нацелена на разработку архитектуры, объединяющей перечисленные типы контекстных сигналов, и закрывает выявленный исследовательский пробел.

\begin{longtable}{p{3.5cm} p{3cm} p{1cm} p{2cm} p{4cm}}
    \caption{Сводная таблица ключевых работ по направлениям}
    \label{tab:ch1:papers} \\
    \toprule
    \textbf{Работа} & \textbf{Авторы} & \textbf{Год} & \textbf{Venue} & \textbf{Направление} \\
    \midrule
    \endfirsthead
    \multicolumn{5}{c}{\tablename\ \thetable{} (продолжение)} \\
    \toprule
    \textbf{Работа} & \textbf{Авторы} & \textbf{Год} & \textbf{Venue} & \textbf{Направление} \\
    \midrule
    \endhead
    \bottomrule
    \endfoot
    GRU4Rec & Hidasi et al. & 2016 & ICLR & Сессионные рек. (RNN) \\
    NARM & Li et al. & 2017 & CIKM & Сессионные рек. (RNN+Attention) \\
    STAMP & Liu et al. & 2018 & KDD & Сессионные рек. (MLP+Attention) \\
    SASRec & Kang \& McAuley & 2018 & ICDM & Последовательные рек. (Transformer) \\
    BERT4Rec & Sun et al. & 2019 & CIKM & Последовательные рек. (Bidirectional) \\
    SR-GNN & Wu et al. & 2019 & AAAI & Графовые сессионные рек. \\
    GC-SAN & Xu et al. & 2019 & IJCAI & GNN + Self-Attention \\
    TiSASRec & Li et al. & 2020 & WSDM & Временные интервалы \\
    GCE-GNN & Wang et al. & 2020 & SIGIR & Глобальный контекст графа \\
    TAGNN & Yu et al. & 2020 & SIGIR & Target-aware GNN \\
    LightGCN & He et al. & 2020 & SIGIR & Графовая кол. фильтрация \\
    Time2Vec & Kazemi et al. & 2020 & ICLR & Периодическое кодирование времени \\
    MEANTIME & Cho et al. & 2020 & RecSys & Смесь временных эмбеддингов \\
    DCN V2 & Wang et al. & 2021 & WWW & Cross-network \\
    S\textsuperscript{3}-Rec & Zhou et al. & 2020 & CIKM & Self-supervised pretraining (attr.) \\
    Transformers4Rec & Moreira et al. & 2021 & RecSys & NLP$\to$RecSys фреймворк \\
    FMLP-Rec & Zhou et al. & 2022 & WWW & FFT-фильтры вместо attention \\
    MB-STR & Yuan et al. & 2022 & SIGIR & Multi-behavior Transformer \\
    HRNN & Quadrana et al. & 2017 & RecSys & Иерархические сессионные рек. (RNN) \\
    HCAN & Cui et al. & 2019 & Neurocomputing & Иерархический трансформер (сессии) \\
    Session Tokens & Seol et al. & 2022 & SIGIR & Сессионные разделители \\
    CARCA & Rashed et al. & 2022 & RecSys & Перекрёстное внимание + атрибуты \\
    DuoRec & Qiu et al. & 2022 & WSDM & Контрастивное обучение (session) \\
    MOJITO & Tran et al. & 2023 & SIGIR & Смесь временных режимов внимания \\
    BSARec & Shin et al. & 2024 & AAAI & Частотный индуктивный сдвиг \\
    PISA & Tran et al. & 2024 & RecSys & Session-уровень + ACT-R память \\
    Query-Aware SR & He et al. & 2022 & CIKM & Поисковые запросы в рек. \\
    HG-GNN & Pang et al. & 2022 & WSDM & Гетерогенный граф \\
    TorchRec & Ivchenko et al. & 2022 & RecSys & Промышленная инфраструктура \\
    gSASRec & Petrov \& Macdonald & 2023 & RecSys & Калибровка потерь \\
    LinRec & Liu et al. & 2023 & SIGIR & Линейное внимание \\
    Embedding Collapse & Guo et al. & 2024 & ICML & Dimensional collapse \\
    HSTU & Zhai et al. & 2024 & ICML & Генеративные рек. трансдьюсеры \\
    UDITSR & --- & 2024 & KDD & Unified Search+Rec \\
    Mamba4Rec & Liu et al. & 2024 & KDD WS & SSM для рек. \\
    TO-RoPE & Wei et al. & 2025 & arXiv & RoPE для рек. \\
    SynerGen & Amazon & 2025 & arXiv & Генеративный unified S+R \\
\end{longtable}
